\section{Joint argument with Chang 2026 structural reduction}

Chang's structural reduction \cite{chang2026reduction} reduces the Collatz
problem to a fixed-modulus, single-bit balance problem: modulus \(32\), bit 4,
and a sparse subsequence of burst-ending times for the dominant
\(n\equiv 1\pmod 8\) class.  In Chang's notation, the Map Balance Theorem 4.2
counts the residues in the burst-to-gap set and proves the exact
\(|C_3(K)-C_7(K)|=1\) imbalance, with perfect balance inside the dominant
\(n\equiv 1\pmod 8\) subclass.  His Open Problem, Eq. 16, asks for every odd
starting value above the stated threshold to keep the empirical split between
\(n_{t_i}\equiv 9\pmod {32}\) and \(n_{t_i}\equiv 25\pmod {32}\) within a
tolerance \(\delta<\delta_{\max}\).  The local compatibility analysis records
this as a compatible reduction, not a closed proof
\(\artifact{chang_2603_25753_compatibility.md}{chang\_2603\_25753\_compatibility.md}\).

The notation correspondence is direct.  Chang's map \(T(n)\) is this paper's
\(\Syr(n)\).  Chang's burst indicator
\[
  X_t=\mathbf{1}[n_t\equiv 1\pmod 4]
\]
is this framework's post-exit indicator.  Chang works at modulus \(32\), with
a modulus \(256\) fiber refinement; the extended Foster audit now reaches
residue quotients modulo \(2^k\) for \(k\in\{2,3,4,5,6,7,8\}\).

The joint argument chain is:
\[
\begin{gathered}
  \text{Foster }m=16,\ \epsilon=0.1\text{ in Section 5}\\
  \Longrightarrow
  \text{Markov ergodicity plus Birkhoff almost-sure equidistribution mod }2^k\\
  \Longrightarrow
  \text{assumption: every integer orbit is in the Markov-measure-1 set}\\
  \Longrightarrow
  \text{every-orbit equidistribution}\\
  \Longrightarrow
  \text{Chang Map Balance Theorem 4.2 plus bit-4 reduction}\\
  \Longrightarrow
  \text{every-orbit bit-4 balance within }\delta_{\max}\\
  \Longrightarrow
  \text{Chang block-TV budget Eq. 3 in the companion paper
  \cite{chang2026collatz}.}
\end{gathered}
\]
This is a joint argument: Foster drift supplies the residue-Markov
equidistribution input, and Chang's compatible reduction converts the relevant
mod \(32\) balance into the block-TV budget.  The conditional step is exactly
where no proof claim is made.

The three technical verifications named in the compatibility analysis are now
empirically closed at finite windows.  First, the \(\delta_{\max}\) analysis
derives \(\delta_{\max}\approx 0.119\) from an \(R(K)\)-weighted allocation
against Chang's Eq. 34
\(\artifact{delta_max_quantitative.md}{delta\_max\_quantitative.md}\).
Second, the Foster condition holds at \(k=7,8\), matching Chang's modulus
\(256\) fiber resolution: \(m=16\) clears the \(\epsilon=0.1\) margin, the
weakest \(k=8\) CI upper bound at \(m=16\) is \(-0.24276696844795248\), and
the \(k\le 6\) cross-check has
\(\texttt{cross\_check\_max\_disagreement}=0.0\)
\(\artifact{m_step_foster_drift_k8.json}{m\_step\_foster\_drift\_k8.json}\).
Third, the direct bit-4 balance audit finds that the Foster plus \(1/\sqrt m\)
envelope holds for every sampled orbit, with deepest-window mean
\(\delta=0.11212329012096864\), just below the \(R(K)\)-weighted
\(\delta_{\max}\approx 0.119\)
\(\artifact{chang_bit4_balance_audit.json}{chang\_bit4\_balance\_audit.json}\).
Thus the joint argument is supported at finite windows.

The residual is exactly the Tao wall.  The verifications closed three
technical gaps, but they did not close the conditional step from
Markov-measure-1 residue-chain behavior to every deterministic integer orbit.
Equivalently, the joint argument reduces Collatz to: every integer orbit lies
in the Markov-measure-1 set on which the residue chain equidistributes.  This
is the Tao 2019 distributional-to-pointwise obstruction
\cite{tao2019almost}, but in a sharper Markov-measure formulation rather than
only a natural-density formulation.

There is one further numerical correspondence worth isolating.  Chang's
phantom-gain sum gives
\[
  \sum_{K=3}^{500} R(K)=0.0882362530512702,
\]
while this framework's renewal Cramér rate is near
\(J_{\mathrm{renewal}}=0.08372911906309355\) in the saved calibration, about
five percent apart.  The direct identity audit recomputes the necklace sum and
finds \(J_{\mathrm{renewal}}=0.08346903426255636\) with bootstrap CI
\([0.08314605994133623,0.0837883621454151]\), outside Chang's sum; the tested
alternate definitions also do not match
\(\artifact{jazz_constant_chang_R_K_identity.json}{jazz\_constant\_chang\_R\_K\_identity.json}\).
The harder \(h_K\)-weighted reconciliation remains future work, so the present
reading is that the five-percent gap is structural in this finite diagnostic.

The empirical comparison is asymmetric in a useful way.  This framework's
Foster condition is verified on \(1000000\) sampled orbits, with
publication-grade confidence intervals and exact cross-checks.  Chang's Route
A reports Lyapunov exponents \(\lambda\in[-0.14,-0.09]\) on five orbits with
eyeball estimates \cite{chang2026reduction}.  The composition strictly
strengthens both: Chang provides the structural one-bit target, while this
framework provides the stronger finite empirical Foster input.  It remains a
finite empirical diagnostic, and the residual is exactly the Tao wall.
